\documentclass[12pt,a4paper]{article}

\usepackage[utf8]{inputenc}
\usepackage[T1]{fontenc}
\usepackage[brazil]{babel}
\usepackage{lmodern}
\usepackage{geometry}
\usepackage{microtype}
\usepackage{array}
\usepackage{tabularx}
\usepackage{enumitem}
\usepackage{xcolor}
\usepackage{fancyhdr}

\geometry{top=1.7cm,bottom=1.7cm,left=1.8cm,right=1.8cm}
\setlength{\parindent}{0pt}
\setlength{\parskip}{0.45em}
\setlist[itemize]{leftmargin=*,itemsep=0.25em,topsep=0.2em}
\definecolor{azul}{HTML}{173B57}
\definecolor{cinza}{HTML}{F2F4F7}

\pagestyle{fancy}
\setlength{\headheight}{15pt}
\fancyhf{}
\lhead{Sequência didática de \textit{machine learning}}
\rhead{Aula 3: SVM linear}
\cfoot{\thepage}
\renewcommand{\headrulewidth}{0.3pt}

\newcommand{\linha}{\rule{\linewidth}{0.35pt}}
\newcommand{\bloco}[1]{%
  \vspace{0.25em}
  \noindent\colorbox{cinza}{\parbox{0.96\linewidth}{#1}}%
  \vspace{0.3em}
}

\begin{document}

\begin{center}
  {\LARGE\bfseries\color{azul} Aula 3: SVM linear}\par
  \vspace{0.25em}
  {\large Fronteira de decisão, margem e vetores de suporte}\par
  \vspace{0.35em}
  {\small Roteiro de aplicação da sequência didática}\par
\end{center}

\bloco{\textbf{Lugar desta aula na pesquisa.} Este é o terceiro encontro da sequência de oito aulas. A atividade apresenta uma forma de classificação supervisionada: exemplos já rotulados ajudam a posicionar uma fronteira entre duas classes. A exploração individual é formativa; o professor registra dificuldades recorrentes da turma.}

\section*{1. Identificação}

\begin{tabularx}{\linewidth}{@{}>{\bfseries}p{4.1cm}X@{}}
Turma: & \rule{8.5cm}{0.35pt} \\
Data: & \rule{8.5cm}{0.35pt} \\
Duração: & 50 minutos \\
Organização: & Exploração individual e síntese oral coletiva. \\
Recursos: & Plataforma no módulo SVM, computadores do laboratório ou Chromebooks, quadro e folha de atividade. \\
\end{tabularx}

\section*{2. Objetivos da aula}

Ao final do encontro, espera-se que os estudantes reconheçam uma fronteira de decisão entre duas classes, expliquem que a margem é o espaço entre a fronteira e os exemplos mais próximos e identifiquem esses exemplos como vetores de suporte.

\textbf{Relação curricular:} EM13CO10, EM13CO11 e EM13CO12.

\textbf{Conceitos em foco:} dado, rótulo, classificação, fronteira de decisão, margem, vetor de suporte e limite do modelo.

\section*{3. Preparação antes da aula}

\begin{itemize}
  \item Abrir o módulo SVM e confirmar que aparecem duas classes de amostras fictícias de sementes.
  \item Conferir se a fronteira, as duas linhas da margem e os vetores de suporte estão legíveis no projetor.
  \item Confirmar que cada estudante terá acesso individual ao módulo, em computador do laboratório ou Chromebook.
  \item Reservar um espaço no quadro para desenhar duas linhas separadoras e comparar suas margens.
\end{itemize}

\bloco{\textbf{Cuidado didático.} As classes A e B e as medidas das sementes são sintéticas. A linha não revela uma propriedade natural das sementes: ela é uma regra construída a partir dos exemplos selecionados.}

\section*{4. Desenvolvimento da aula}

\renewcommand{\arraystretch}{1.18}
\begin{tabularx}{\linewidth}{|>{\centering\arraybackslash}p{1.65cm}|p{5.0cm}|X|}
\hline
\textbf{Tempo} & \textbf{Ação com a turma} & \textbf{Mediação docente} \\
\hline
0--7 min & Retomar a regressão linear e contrastar previsão numérica com escolha entre duas classes. & Perguntar: ``Na aula anterior, a saída era uma distância. Se agora a saída for A ou B, qual tipo de problema temos?'' \\
\hline
7--12 min & Apresentar as sementes fictícias por comprimento e largura. Marcar A e B no plano e desenhar duas possíveis linhas de separação. & Evitar dizer que uma linha é automaticamente melhor. Perguntar qual linha deixa mais espaço até os pontos próximos. \\
\hline
12--19 min & Pedir que cada estudante preveja o que acontece com a margem quando uma fronteira se aproxima da classe A. & Solicitar uma justificativa com os pontos mais próximos da linha, antes de abrir o controle. \\
\hline
19--34 min & Explorar o módulo de forma guiada. Posicionar a fronteira e depois comparar com a sugestão do algoritmo. & Nomear fronteira, margem e vetores de suporte a partir do que está visível. Manter a atenção em uma alteração de cada vez. \\
\hline
34--43 min & Incluir o ponto próximo e comparar o tamanho da margem. Mover o ponto novo para observar em que lado da linha ele fica. & Perguntar como a entrada do novo ponto é ligada à classe pela posição em relação à fronteira. Retomar que a margem pode diminuir. \\
\hline
43--50 min & Realizar a atividade individual breve e encerrar com uma síntese oral. & Esclarecer o enunciado sem indicar a resposta. Retomar que a regra depende dos dados, dos rótulos e da forma de separação escolhida. \\
\hline
\end{tabularx}

\section*{5. Perguntas para conduzir a conversa}

\begin{itemize}
  \item ``O que a linha separa neste gráfico?''
  \item ``Quais exemplos estão mais próximos da fronteira?''
  \item ``O que acontece com a margem quando a fronteira se aproxima de um ponto?''
  \item ``Em que lado da fronteira o ponto novo ficou?''
  \item ``Uma separação boa neste conjunto garante que funcionará em qualquer conjunto?''
\end{itemize}

\section*{6. Evidência a observar}

Observar se os estudantes distinguem dados, rótulos e fronteira; se reconhecem que os vetores de suporte são os exemplos mais próximos da linha; e se conseguem justificar a classificação de um ponto novo pelo lado em que ele se encontra. No registro docente, anotar confusões entre margem e distância entre quaisquer pontos, justificativas formuladas individualmente e condições de uso da plataforma que tenham afetado a aula.

\newpage

\begin{center}
  {\LARGE\bfseries\color{azul} Atividade individual breve}\par
  \vspace{0.2em}
  {\large Aula 3: fronteira, margem e classificação}\par
  \vspace{0.35em}
  {\small Material formativo para o final da aula}\par
\end{center}

\bloco{\textbf{Orientação ao estudante.} Esta atividade não vale nota. Responda individualmente e explique como você pensou. Use o que foi observado na plataforma e o que foi discutido na aula.}

\textbf{Código do participante:} \rule{5cm}{0.4pt}
\hfill
\textbf{Data:} \rule{3cm}{0.4pt}

\section*{Situação}

Em um gráfico, exemplos da classe A ficam no lado inferior esquerdo e exemplos da classe B ficam no lado superior direito. Uma linha foi desenhada para separar as duas classes. Os pontos A4 e B1 são os mais próximos da linha.

\section*{Questão 1}

Por que A4 e B1 podem ser chamados de vetores de suporte? Explique a relação deles com a margem.

\vspace{0.25em}
\linha\par\vspace{0.5em}
\linha\par\vspace{0.5em}
\linha

\section*{Questão 2}

Imagine que a linha seja deslocada para mais perto de A4, mantendo a mesma inclinação. O que tende a acontecer com a margem? Explique.

\vspace{0.25em}
\linha\par\vspace{0.5em}
\linha\par\vspace{0.5em}
\linha

\section*{Questão 3}

Um novo ponto fica no mesmo lado da fronteira que os exemplos B. Qual classe a SVM indicará para ele? O que poderia dificultar essa classificação em outro conjunto de dados?

\vspace{0.25em}
\linha\par\vspace{0.5em}
\linha\par\vspace{0.5em}
\linha

\end{document}
